\section{Certainty ratings}\label{app:C}

For each finding in \cref{tab:keyfindings} we rated the certainty of the evidence in the spirit of GRADE \cite{guyatt2008grade}. GRADE starts non-randomised evidence at low, and benchmark evaluations are non-randomised. The rules below are our adaptation for single-model benchmark studies, not part of GRADE.

\paragraph{Levels.} A finding is very low if it rests on a single study, on studies none of which is of stronger design, or, when it concerns hosted \jev{} or typed models generally, on open-model evidence only. It is low if at least two studies agree in direction and at least one is of stronger design. It is moderate if at least three stronger-design studies, independent of one another, agree in direction and at least one of them is preregistered or independently replicated. A finding with several clauses takes the lowest rating of its clauses.

\paragraph{Counting.} Studies sharing an author count as one study. A study counts as an independent replication only if its authors differ from the original's, it uses its own data or a reimplementation rather than re-running the original code on the original data, and its design is at least moderate. Each finding has a scope: hosted \jev{}, open models, or typed models generally. Evidence is counted per system arm, so a study that tested hosted \jev{} and open models has a hosted arm and an open arm. A finding about hosted \jev{} or about typed models generally is rated on hosted-model evidence alone. Open-model evidence can corroborate such a finding but cannot raise, lower or contest its rating. Open-model evidence that points against such a finding is recorded with the direction contradict-open, so that it remains visible in the data and in \cref{tab:keyfindings} (``open, against''), but it does not enter the rating.

\paragraph{Opposing evidence.} A study opposes a finding when it reports, for the system and setting the finding names, a result on the outcome the finding states that runs the other way. A study with results in both directions counts on the side of its primary result. A result for another system, another outcome or an excluded setting neither supports nor opposes. Studies are ranked by design rating (stronger above moderate above weaker) and, within a rating, preregistered or replicated studies above the others. Any opposing study marks the finding contested. If the strongest opposing study ranks at least as high as the strongest supporting one, the rating also falls one level. For findings about escalation, a second stage counts as a stronger reasoner when, in the same study, it was more accurate than the first stage on the task on which escalation was tested; a difference within sampling error, or accuracy measures that point in opposite directions, makes it a peer.

\paragraph{Application.} The rules were applied to every finding with all coded studies that bear on it, not only those cited for it. Every coded study's assignment to each finding, with a reason, is in \nolinkurl{data/certainty_evidence.csv}, and \nolinkurl{scripts/rate_certainty.py} computes the ratings from it and checks the table of the paper. The last rated row of \cref{tab:keyfindings} is a tally of comparator codes rather than a set of studies; the rules do not apply to it, and we rate it low by judgement. The comparator field was second-coded on 148 pairs ($\kappa = 0.72$) but coded once on the others, and the tally counts studies rather than measuring a difference, so it cannot rate higher than the findings it summarises. \Cref{tab:C1} lists, for every finding, the studies counted for and against it.

% Generated by scripts/rate_certainty.py --derivation (do not edit the table by hand)
% Generated by scripts/rate_certainty.py --derivation; do not edit by hand.
\begingroup\scriptsize\setlength{\tabcolsep}{3pt}
\begin{longtable}{@{}>{\raggedright\arraybackslash}p{0.04\linewidth}>{\raggedright\arraybackslash}p{0.25\linewidth}>{\raggedright\arraybackslash}p{0.36\linewidth}>{\raggedright\arraybackslash}p{0.17\linewidth}>{\raggedright\arraybackslash}p{0.10\linewidth}@{}}
\caption{Derivation of each certainty rating. For every finding (or clause of a finding), the studies counted as supporting and opposing it on the counted system arm (hosted, unless the finding concerns open models), with their design (S stronger, M moderate, W weaker; * preregistered or independently replicated). Studies sharing an author count once. Rules in \oa{C}; a finding with several clauses takes its lowest clause. Source: \nolinkurl{data/certainty_evidence.csv}.}\label{tab:C1}\\
\toprule
Id & Finding (scope) & Supporting studies & Opposing studies & Rating \\
\midrule
\endfirsthead
\toprule
Id & Finding (scope) & Supporting studies & Opposing studies & Rating \\
\midrule
\endhead
\bottomrule
\endfoot
1a & ranks harmful or misaligned content well zero-shot (hosted) & \cite{arx01079}\,M, \cite{arx03324}\,S, \cite{arx03935}\,M, \cite{arx04985}\,M, \cite{arx07953}\,S, \cite{arx29429}\,S, \cite{arx33401}\,S, \cite{jkf87_2026replication}\,W, \cite{zen22885291}\,S, \cite{zen23075657}\,S & -- & low \\
\addlinespace
1b & at a fraction of the cost (hosted) & \cite{arx24052}\,S, \cite{arx24574}\,S*, \cite{arx27678}\,S, \cite{arx29769}\,S, \cite{zen22849329}\,S & -- & moderate \\
\addlinespace
1c & often less accurate than the best LLMs (hosted) & \cite{arx01079}\,M, \cite{arx02293}\,S, \cite{arx03324}\,S, \cite{arx03935}\,M, \cite{arx08775}\,M, \cite{arx08829}\,W, \cite{arx24574}\,S*, \cite{arx27678}\,S, \cite{arx33401}\,S, \cite{arx34963}\,W, \cite{zen22849329}\,S, \cite{zen23005202}\,W & -- & moderate \\
\addlinespace
2 & Fixed thresholds do not transfer; local fitting is needed (both) & \cite{arx00346}\,S, \cite{arx01006}\,S, \cite{arx02046}\,W, \cite{arx02267}\,S, \cite{arx03387}\,S, \cite{arx04985}\,M, \cite{arx24574}\,S*, \cite{arx26550}\,S, \cite{arx27607}\,M, \cite{arx29429}\,S, \cite{arx39496}\,S, \cite{jkf87_2026replication}\,W, \cite{zen22935043}\,S, \cite{zen22952571}\,S, \cite{zen23050542}\,M & \cite{arx33689}\,S* & low, contested \\
\addlinespace
3 & Errors concentrate in sub-types hidden by aggregates (hosted) & \cite{alphaRAG}\,W, \cite{arx00213}\,S, \cite{arx02046}\,W, \cite{arx02267}\,S, \cite{arx03324}\,S, \cite{arx07953}\,S, \cite{arx08675}\,S, \cite{arx10321}\,S, \cite{arx33401}\,S, \cite{zen22846597}\,W & -- & low \\
\addlinespace
4 & Inserted opinions, optimised context and lexical lures steer decisions (both) & \cite{arx01006}\,S, \cite{arx01834}\,S, \cite{arx02048}\,M*, \cite{arx04985}\,M, \cite{arx30243}\,W, \cite{arx31142}\,S & -- & low \\
\addlinespace
5 & Instruction injection rarely selects the attacker target (hosted) & \cite{arx28613}\,S* & \cite{arx04985}\,M, \cite{arx31142}\,S, \cite{checkpoint2026jev}\,W & very low, contested \\
\addlinespace
6a & option names can override definitions (both) & \cite{arx00346}\,S, \cite{arx07953}\,S, \cite{arx26758}\,S & -- & low \\
\addlinespace
6b & neutral identifiers remove most of it (both) & \cite{arx00346}\,S, \cite{arx07953}\,S, \cite{arx26758}\,S & -- & low \\
\addlinespace
7 & Defaults to one culture's values (hosted) & \cite{arx36399}\,S & -- & very low \\
\addlinespace
8a & well calibrated on familiar closed-choice tasks (hosted) & \cite{arx00346}\,S, \cite{arx01006}\,S, \cite{arx03935}\,M, \cite{arx05107}\,M, \cite{arx06625}\,S, \cite{arx09188}\,S, \cite{arx09896}\,M, \cite{arx34024}\,W, \cite{arx37647}\,W, \cite{qin2026zhdecisionbench}\,W, \cite{zen22885291}\,S & -- & low \\
\addlinespace
8b & mixed against frontier models (hosted) & \cite{arx02293}\,S, \cite{arx06625}\,S, \cite{arx24052}\,S, \cite{arx24574}\,S*, \cite{arx26550}\,S, \cite{arx34024}\,W, \cite{zen22885291}\,S & -- & moderate \\
\addlinespace
8c & calibration is local (hosted) & \cite{alphaRAG}\,W, \cite{arx00213}\,S, \cite{arx00346}\,S, \cite{arx01006}\,S, \cite{arx02267}\,S, \cite{arx07953}\,S, \cite{arx09937}\,S, \cite{arx10321}\,S, \cite{arx24574}\,S*, \cite{arx27607}\,M, \cite{arx29429}\,S, \cite{arx29769}\,S, \cite{arx35342}\,S, \cite{arx36154}\,S, \cite{arx37470}\,W, \cite{grazian2026calibrated}\,S, \cite{zen22846597}\,W, \cite{zen22935043}\,S, \cite{zen22945279}\,W, \cite{zen22980293}\,S*, \cite{zen23032384}\,S*, \cite{zen23088660}\,S, \cite{zen23179064}\,S*, \cite{zen23221456}\,S & -- & moderate \\
\addlinespace
9a & 'I don't know' rarely chosen when correct (hosted) & \cite{arx01006}\,S, \cite{arx07730}\,M, \cite{arx34024}\,W, \cite{zen23119392}\,W & -- & low \\
\addlinespace
9b & 'none' chosen inconsistently (hosted) & \cite{arx01006}\,S, \cite{arx03387}\,S, \cite{arx06425}\,W, \cite{arx39496}\,S, \cite{zen22864020}\,W & -- & low \\
\addlinespace
10a & sufficiency question detects missing information (hosted) & \cite{arx01006}\,S, \cite{arx05107}\,M & -- & low \\
\addlinespace
10b & filters less well than confidence (hosted) & \cite{arx01006}\,S & -- & very low \\
\addlinespace
11 & Escalation to a stronger reasoner saves cost without loss of accuracy (both) & \cite{arx02046}\,W, \cite{arx24574}\,S*, \cite{arx26532}\,M, \cite{arx26550}\,S & \cite{arx02267}\,S & low, contested \\
\addlinespace
12a & acceptance on rule agreement lost no macro-F1 (both) & \cite{arx02048}\,M* & -- & very low \\
\addlinespace
12b & fewer disguised attacks flagged (both) & \cite{arx02048}\,M* & -- & very low \\
\addlinespace
13 & Peer evaluators repeat Jev's most confident errors (hosted) & \cite{arx29769}\,S, \cite{arx33401}\,S & -- & low \\
\addlinespace
14 & Thresholds set in advance miss error targets on new data (both) & \cite{arx02267}\,S, \cite{arx03387}\,S, \cite{arx24574}\,S*, \cite{arx33401}\,S, \cite{zen23050542}\,M & \cite{arx33689}\,S*, \cite{zen22935043}\,S, \cite{zen23075657}\,S & low, contested \\
\addlinespace
15 & Typed models not generally worse than generative evaluators (descriptive tally) & \multicolumn{2}{l}{descriptive tally of comparator codes; rated by judgement} & low \\
\end{longtable}
\endgroup

